\documentclass{article}
\usepackage[T2A]{fontenc}
\usepackage[utf8]{inputenc}
\usepackage{amsthm}
\usepackage{amsmath}
\usepackage{amssymb}
\usepackage{amsfonts}
\usepackage{mathrsfs}
\usepackage[12pt]{extsizes}
\usepackage{fancyvrb}
\usepackage{indentfirst}
\usepackage[
  left=2cm, right=2cm, top=2cm, bottom=2cm, headsep=0.2cm, footskip=0.6cm, bindingoffset=0cm
]{geometry}
\usepackage[english,russian]{babel}

\begin{document}
\section*{Вариант 5}

Производя подстановку для одновременного применения обоих сдвигов, получим:

\[
\begin{cases}
A = 3\delta - b_s = -b - 3a\Delta + 3\delta - 6\Delta^2, \\[4pt]
\begin{aligned}
B &= 3\delta^2 - 2b_s\delta + a_s c_s - 4d_s \\
  &= 3\delta^2 - 2(b + 3a\Delta + 6\Delta^2)\delta + (a + 4\Delta)(c + 2b\Delta + 3a\Delta^2 + 4\Delta^3) \\
  &\quad - 4(d + c\Delta + b\Delta^2 + a\Delta^3 + \Delta^4),
\end{aligned} \\[4pt]
\begin{aligned}
C &= \delta^3 - b_s\delta^2 + (a_s c_s - 4d_s)\delta - a_s^2 d_s - c_s^2 + 4b_s d_s \\
  &= \delta^3 - (b + 3a\Delta + 6\Delta^2)\delta^2 \\
  &\quad + \bigl((a + 4\Delta)(c + 2b\Delta + 3a\Delta^2 + 4\Delta^3) \\
  &\qquad - 4(d + c\Delta + b\Delta^2 + a\Delta^3 + \Delta^4)\bigr)\delta \\
  &\quad - (a + 4\Delta)^2(d + c\Delta + b\Delta^2 + a\Delta^3 + \Delta^4) \\
  &\quad - (c + 2b\Delta + 3a\Delta^2 + 4\Delta^3)^2 \\
  &\quad + 4(b + 3a\Delta + 6\Delta^2)(d + c\Delta + b\Delta^2 + a\Delta^3 + \Delta^4).
\end{aligned}
\end{cases}
\]

Расчёты показывают, что когда \(\delta = \alpha a_s^2 + \beta b_s\), где \(8\alpha + 3\beta = 1\), коэффициенты резольвенты не меняются при любом сдвиге \(x\).

\end{document} 
